\section{Introduction}

\kleos{} is an AI operating system for computer science students (\url{https://kleos-cs.vercel.app}). Beneath its planning, memory and research features sits a layer of small decisions that a language model makes many times a day: which of several competing items deserves attention first, which of two contradictory records to trust, whether a request is specific enough to act on, which items belong in a briefing, and which sit in a workspace the user has not opened. These decisions are frequent and consequential, and a general-purpose model prompted with instructions does not reliably make them the way the product requires (the orchestrated base models score 0.4469 to 0.4755 on the benchmark, Section~\ref{sec:results}).

% Figure 1: system overview. Labels only; counts from docs/experiments.md (Status),
% Kleos-Training-Data README (releases) and docs/evaluation.md (benchmark).
\begin{figure}[t]
\centering
\begin{tikzpicture}[
  node distance=4mm and 6mm,
  box/.style={draw, rounded corners=2pt, align=center, font=\footnotesize, inner sep=4pt, text width=30mm, minimum height=11mm},
  wide/.style={box, text width=38mm},
  arr/.style={-{Latex[length=2mm]}, thick}
]
% Row 1: data
\node[box] (cat) {Scenario catalog\\ 18 scenario families,\\ 12 decision policies};
\node[box, right=of cat] (gen) {Generation\\ 1,350 examples; every\\ answer computed by a policy};
\node[box, right=of gen] (gate) {Promotion gates (14)\\ and sealing\\ (content hash)};
\node[wide, right=of gate] (split) {Group split (sha256, seed 42)\\ train 820 / validation 181 /\\ test 349, JSON format held out};
\draw[arr] (cat) -- (gen);
\draw[arr] (gen) -- (gate);
\draw[arr] (gate) -- (split);
% Row 2: models and evaluation
\node[box, below=10mm of cat] (qlora) {QLoRA fine-tuning\\ NF4, $r=16$, T4 GPUs\\ (Colab, Kaggle)};
\node[box, right=of qlora] (arms) {Arms\\ base + orchestration (arm1)\\ vs fine-tuned (arm2)};
\node[box, right=of arms] (bench) {Benchmark, 349 items\\ 271 answerable / 78\\ should-decline; 78 groups};
\node[wide, right=of bench] (verdict) {Pre-registered verdict\\ paired bootstrap, cluster-\\ resampled by group from H8 on};
\draw[arr] (qlora) -- (arms);
\draw[arr] (arms) -- (bench);
\draw[arr] (bench) -- (verdict);
\draw[arr] (split.south west) ++(8mm,0) |- ++(0,-3mm) -| (qlora.north);
% Row 3: serving
\node[box, below=10mm of arms] (pkg) {Deployment package\\ pinned base revision,\\ adapter SHA-256};
\node[box, right=of pkg] (space) {ZeroGPU Space\\ greedy decoding,\\ reproduction check};
\node[box, right=of space] (app) {\kleos{} application\\ Hermes (fast),\\ Logos (shows its reasoning)};
\draw[arr] (verdict.south) |- ++(0,-3mm) -| (pkg.north);
\draw[arr] (pkg) -- (space);
\draw[arr] (space) -- (app);
\end{tikzpicture}
\caption{The pipeline. The data pipeline (top row) lives in the companion repository and seals each release with a content hash; the training, evaluation and serving code (rows two and three) live in the KLEOS Models repository. The test split is the sealed benchmark; it was never used for selection.}
\label{fig:pipeline}
\end{figure}


\kleos{} therefore defines its decisions as an explicit policy. Each candidate item is scored as the product of a deadline term, an evidence weight and an impact weight; items whose scores fall within a relative margin are declined rather than guessed; an explicit statement by the user outranks an inferred record unless it is stale; items outside the active workspace are flagged rather than acted on; and an ambiguous request is answered with a question. The policy is deterministic, so the correct answer to any scenario is computable. That makes two things possible: a synthetic training set in which every answer is correct by construction, and a benchmark on which a model's judgment is measured against the policy rather than against human raters.

This paper asks whether a small open-weight model can learn that policy from such data under conditions a student can afford: free Colab and Kaggle T4 GPUs, QLoRA adapters rather than full fine-tuning, and no proprietary models. It reports two production models. \hermes{} is a QLoRA adapter on Mistral-Nemo-Instruct 12B. \logos{} is an adapter on Ministral 3 14B, in a first version on the Instruct checkpoint and a second on the Reasoning checkpoint, trained on reasoning traces derived from the same policy. Every hypothesis that was run was written down with its metric and decision rule before the run, every deviation from the protocol is dated and logged, and negative and inconclusive results are reported with the same precision as positive ones. The design follows the pre-registration literature \citep{nosek2018preregistration,vanmiltenburg2021preregistering,bertinetto2021preface} and recent guidance on interval estimates for language-model evaluations \citep{miller2024errorbars}.

% source: EXP:167 (pilot), HER:259-275 (Hermes), EXP:545-553 (H8b), EXP:748-770 (H9).
Three findings follow. First, fine-tuning on 820 synthetic examples moved the benchmark score of Hermes from 0.4755 to 0.8051 (difference $+0.3295$, 95\% CI 0.3068 to 0.3521, every one of the seven tasks $p<0.001$), replicating a pilot on Ministral-8B (0.4744 to 0.8015, seven of seven tasks). Second, moving to the stronger Ministral 3 base did not by itself produce a measurably better model: Logos v0.0.1 differed from Hermes by $-0.0168$ on the answerable items (cluster 95\% CI $-0.0546$ to $+0.0177$), an inconclusive result under the pre-registered rule. Third, Logos v0.0.2, trained on the data release that adds policy-derived reasoning traces, was measurably better than Hermes on the answerable items ($+0.0409$, cluster 95\% CI $+0.0040$ to $+0.0780$, $p=0.031$); its gain on the should-decline items did not reach significance by groups. Each model is a single training run, so these intervals cover evaluation noise only, and the v0.0.2 comparison confounds three changes made together.

The contributions are:
\begin{itemize}
  \item A sealed synthetic dataset and benchmark for a product decision policy: 1,350 examples over 18 scenario families and 12 policies, identified by content hashes and rebuilt byte for byte by the public pipeline, whose 349-item test split is byte-identical across the two releases used here (Section~\ref{sec:data}).
  \item A pre-registered protocol for fine-tuning experiments on free compute: dated hypotheses H1 to H9, six of them with decision rules, a cluster-bootstrap amendment adopted before the second model line was trained, and a log of thirteen deviations, D1 to D13, each with its effect on the result (Section~\ref{sec:protocol}, Appendices~\ref{app:deviations} and \ref{app:artifacts}).
  \item The H1 result and its replication across two base models, with sub-score, consistency and abstention analyses (Sections~\ref{sec:results} and \ref{sec:analysis}).
  \item The H8 and H9 comparisons between model lines, with their confounds declared before the runs (Sections~\ref{sec:results} and \ref{sec:analysis}).
  \item A serving path on free ZeroGPU Spaces with pinned base revisions and adapter hashes, verified by byte-for-byte reproduction of the frozen evaluation outputs: 9 of 9 for Hermes and 8 of 9 for Logos (Section~\ref{sec:serving}).
\end{itemize}

Figure~\ref{fig:pipeline} shows how the data pipeline, the training and evaluation harness and the serving path fit together, and Section~\ref{sec:related} places the work among its neighbours. Throughout, every number is copied from a frozen record in the public repository, and each table names the record it reproduces. Analyses are labelled pre-registered, with their hypothesis identifier, or exploratory. Section~\ref{sec:limitations} collects the limitations and Section~\ref{sec:repro} states what can be reproduced, and by whom.
